% ═══════════════════════════════════════════════════════════════════
%  Section 5 — Results and Analysis
% ═══════════════════════════════════════════════════════════════════

\label{sec:results}

This section reports results for the six completed experiments (A--E, H) described in Section~\ref{sec:experiments}.  We begin with infrastructure-level observations that apply across all experiments, then present per-experiment findings, and conclude with a cross-experiment synthesis.

% ───────────────────────────────────────────────────────────────────
\subsection{LLM Quality and Gate Summary}
\label{sec:results-quality}

Before examining experimental outcomes, we report the reliability of the LLM backend and the feasibility gate's involvement across all 2,770 sessions.

\begin{table}[ht]
\centering
\caption{LLM backend quality and feasibility gate involvement across all experiments.}
\label{tab:llm-quality}
\begin{tabular}{@{}lr@{}}
\toprule
Metric & Value \\
\midrule
LLM fallback rate (format failures) & 0.0\% \\
Post-LLM override rate & 0.0\% \\
Gate-driven acceptances (fraction of all turns) & $\sim$40\% \\
LLM-driven decisions (fraction of all turns) & $\sim$60\% \\
\bottomrule
\end{tabular}
\end{table}

% TODO: Compute exact gate percentage from JSONL logs if more precision is needed.

The \texttt{llama3.2:3b} model produced valid JSON on every call across all experiments, requiring zero fallback actions.  The post-LLM feasibility check (which overrides \textsc{counter} or \textsc{reject} to \textsc{accept} when the standing offer is feasible) was never triggered, indicating that when the gate did not pre-empt the LLM call, the model never rejected a clearly profitable offer.

The gate drove approximately 40\% of all acceptances.  In fixed-scenario experiments (A--C), the gate fires deterministically after the seller's round-1 counter-offer whenever that counter is within the buyer's constraints, settling all sessions at round~2.  In market-mode experiments (D--E), gate involvement varies per session depending on the randomly drawn ZOPA for each pair.

% TODO: Break down gate involvement rate per experiment if per-experiment JSONL analysis is performed.

Experiment~E produced 2 JSON parse failures and 1 LLM timeout across 900 sessions; the retry and fallback mechanisms handled all three without data loss.

% ───────────────────────────────────────────────────────────────────
\subsection{Experiment A: Concession Dynamics}
\label{sec:results-concession}

Table~\ref{tab:concession} reports the mean offer price at each observed round for the three agent-type conditions.  All 450 sessions (150 per condition) resulted in deals, confirming H-A1.

\begin{table}[ht]
\centering
\caption{Mean offer price by agent type and round ($n = 150$ per condition, pooled across 3 seeds).  Concession $\Delta$ is the absolute difference between the seller's round-1 counter and the buyer's round-0 opening.}
\label{tab:concession}
\begin{tabular}{@{}lccc@{}}
\toprule
Condition & Round 0 (Buyer opens) & Round 1 (Seller counters) & Concession $\Delta$ \\
\midrule
\texttt{rule\_based}       & \$60.00 $\pm$ \$0.00 & \$101.33 $\pm$ \$0.00 & \$41.33 \\
\texttt{llm\_reactive}     & \$65.46 $\pm$ \$0.90 & \$103.78 $\pm$ \$3.05 & \$38.32 \\
\texttt{llm\_deliberative} & \$66.74 $\pm$ \$0.67 & \$96.57 $\pm$ \$2.72  & \$29.83 \\
\bottomrule
\end{tabular}
\end{table}

All sessions settled at round~2 via the feasibility gate, which fired after the seller's round-1 counter-offer.  This limits the concession data to two observations per session.

\paragraph{Finding 1: LLM agents open less aggressively than the rule-based baseline.}
Both LLM conditions open \$5--7 higher than the rule-based agent (\$65--67 vs.\ \$60).  The rule-based buyer opens at exactly 50\% of its value ceiling, while the LLM agents do not reproduce this heuristic precisely.

\paragraph{Finding 2: Chain-of-thought prompting produces more conservative counter-offers.}
The deliberative agent's seller counter (\$96.57) is \$7.21 lower than the reactive agent's (\$103.78), yielding a concession gap of \$29.83 versus \$38.32.  This is a prompting effect---the BELIEFS$\to$TARGET$\to$STRATEGY$\to$ACTION chain leads the model to hold a firmer position---rather than evidence of Rubinstein-style time-preference patience.  H-A3 is confirmed: structured reasoning alters counter-offer behaviour.

\paragraph{Finding 3: Rule-based agent produces zero variance.}
Standard deviation of \$0.00 across all 150 sessions per round confirms that the seeded RNG and linear concession formula are fully deterministic, validating the platform's reproducibility guarantees.

\paragraph{Limitation.}
The feasibility gate forces settlement at round~2 in all sessions, meaning only two data points per session are observed.  A full concession trajectory over 10 rounds is not available in this experiment.  Future work should test with the gate disabled or a narrower ZOPA to observe multi-round concession curves.

% ───────────────────────────────────────────────────────────────────
\subsection{Experiment B: Anchoring Effects}
\label{sec:results-anchoring}

Table~\ref{tab:anchoring} reports deal rates, first offers, and settlement prices across the three ZOPA conditions.

\begin{table}[ht]
\centering
\caption{Anchoring experiment outcomes by condition ($n = 150$ per condition, pooled across 3 seeds).  Seller cost is fixed at \$80 in all conditions.}
\label{tab:anchoring}
\begin{tabular}{@{}lcccc@{}}
\toprule
Condition & Deal rate & First offer & Final price & Buyer surplus \\
\midrule
\texttt{anchor\_low} (ZOPA = \$20)  & 2.7\% (4/150) & \$55.00 $\pm$ \$0.00 & \$96.40 $\pm$ \$0.00 & \$3.60 \\
\texttt{anchor\_mid} (ZOPA = \$40)  & 100\% (150/150) & \$65.51 $\pm$ \$1.03 & \$103.48 $\pm$ \$3.25 & \$16.52 \\
\texttt{anchor\_high} (ZOPA = \$70) & 100\% (150/150) & \$75.00 $\pm$ \$0.00 & \$101.55 $\pm$ \$5.50 & \$48.45 \\
\bottomrule
\end{tabular}
\end{table}

First-offer / final-price Pearson $r = -0.148$ across all deals.

\paragraph{Finding 4: ZOPA width is the primary determinant of deal success.}
The \texttt{anchor\_low} condition collapses to a 2.7\% deal rate compared to 100\% in the other two conditions.  With ZOPA~=~\$20, the seller's opening counter of approximately \$104 exceeds the buyer's hard cap of \$100, making the offer infeasible.  The feasibility gate cannot fire, and the LLM must negotiate the seller below its natural opening position---a task it nearly always fails.  H-B1 is confirmed.

\paragraph{Finding 5: Classical anchoring bias is absent (null result).}
Pearson $r = -0.148$ indicates no meaningful positive relationship between first offer and final price.  The \texttt{anchor\_high} condition (first offer \$75) does not produce higher settlement prices than \texttt{anchor\_mid} (first offer \$65); final prices converge to the \$101--103 range regardless.  H-B2 is rejected; H-B2$'$ is supported.

This is a departure from human bargaining results, where first offers reliably anchor final outcomes~\cite{galinsky2001}.  One possible explanation is that the accept-condition block in the prompt provides explicit surplus information, reducing the model's reliance on the first offer as an informational signal.  This hypothesis has not been tested via ablation.

\paragraph{Design confound.}
The independent variable simultaneously manipulates ZOPA width and first-offer level, making it impossible to fully isolate anchoring from ZOPA feasibility effects.  The Pearson $r$ metric partially addresses this by measuring within-deal correlation, but the confound remains.

% ───────────────────────────────────────────────────────────────────
\subsection{Experiment C: Deadline Pressure}
\label{sec:results-deadline}

Table~\ref{tab:deadline} reports settlement timing across the three horizon conditions.

\begin{table}[ht]
\centering
\caption{Deadline experiment outcomes by time horizon ($n = 150$ per condition, pooled across 3 seeds).}
\label{tab:deadline}
\begin{tabular}{@{}lcccc@{}}
\toprule
Max rounds & Deal rate & Avg closing round & SD & Last-2-round share \\
\midrule
4  & 100\% & 3.00 & 0.00 & 100.0\% \\
8  & 100\% & 3.00 & 0.00 & 0.0\% \\
16 & 100\% & 3.00 & 0.00 & 0.0\% \\
\bottomrule
\end{tabular}
\end{table}

\paragraph{Finding 6: Deadline horizon does not affect settlement timing (null result).}
All conditions close at exactly round~3 with zero variance, regardless of whether the time limit is 4, 8, or 16 rounds.  The feasibility gate fires deterministically after the seller's round-1 counter-offer, producing a mechanical settlement at round~2 that is fully independent of deadline proximity.  H-C1 is rejected; H-C1$'$ is confirmed.

\paragraph{Finding 7: The \texttt{max\_rounds=4} last-2-round share is an arithmetic artifact.}
The 100\% last-2-round share under \texttt{max\_rounds=4} occurs because round~3 of~4 falls within the ``last 2 rounds'' window.  This is not evidence of deadline-driven concession; it is a consequence of the gate's fixed firing time relative to a short horizon.

\paragraph{Confounds.}
Two design choices compound to produce this null result.  First, the feasibility gate settles deals before deadline proximity becomes relevant.  Second, the prompt steers agents to prefer \textsc{counter} over \textsc{reject}, removing the primary mechanism by which agents could delay settlement.  To study genuine deadline effects, future work should disable the gate and remove the prefer-counter instruction, accepting lower deal rates as part of the measurement.

% ───────────────────────────────────────────────────────────────────
\subsection{Experiment D: Market Price Discovery}
\label{sec:results-market}

Tables~\ref{tab:market-agg} and~\ref{tab:market-early-late} report aggregate market statistics and temporal comparison for the 30-tick market simulation.

\begin{table}[ht]
\centering
\caption{Market dynamics aggregate statistics (30 ticks, single seed).}
\label{tab:market-agg}
\begin{tabular}{@{}lr@{}}
\toprule
Metric & Value \\
\midrule
Mean transaction price & \$89.44 \\
Price range & \$78.12 -- \$110.03 \\
Within-tick price std (mean) & \$20.02 \\
Average liquidity & 52.2\% \\
Liquidity range & 13.3\% -- 80.0\% \\
Price trend slope & +\$0.057 / tick \\
Buyer surplus (mean) & \$32.95 \\
Seller surplus (mean) & \$16.40 \\
\bottomrule
\end{tabular}
\end{table}

\begin{table}[ht]
\centering
\caption{Early vs.\ late market comparison (single seed).}
\label{tab:market-early-late}
\begin{tabular}{@{}lcc@{}}
\toprule
Period & Mean price & Mean liquidity \\
\midrule
Ticks 0--4 (early)   & \$84.18 & 54.7\% \\
Ticks 25--29 (late)  & \$86.57 & 61.3\% \\
\bottomrule
\end{tabular}
\end{table}

% TODO: Generate Figure — price time series (mean ± std band over 30 ticks).
% \begin{figure}[ht]
% \centering
% \includegraphics[width=0.8\textwidth]{figures/market_price_timeseries.pdf}
% \caption{Mean transaction price per tick with within-tick standard deviation band (30 ticks, single seed).}
% \label{fig:market-price}
% \end{figure}

\paragraph{Finding 8: Market prices stabilise rapidly.}
The price trend slope of +\$0.057/tick is economically negligible (total drift $<$~\$2 over 30~ticks against a mean of \$89).  The market stabilises within the first few ticks, consistent with the prediction that decentralised bilateral bargaining converges to a stable price level even without a central auctioneer~\cite{rubinstein1985}.  H-D1 is confirmed.

\paragraph{Finding 9: Law of One Price fails; price dispersion persists.}
Within-tick standard deviation of \$20 means that simultaneous transactions occur at substantially different prices within the same market period.  This is the expected outcome under bilateral bargaining, where each pair negotiates independently and prices reflect the specific ZOPA of each pair rather than a single market-clearing price.  H-D2 is confirmed.

\paragraph{Finding 10: Buyers capture more surplus than sellers.}
Mean buyer surplus (\$32.95) exceeds mean seller surplus (\$16.40) by a factor of two.  The mean transaction price of \$89.44 lies below the ZOPA midpoint for most pairs, consistent with a first-mover advantage: buyers open in round~0, anchoring the negotiation toward the lower end of the ZOPA.  H-D3 is confirmed.

\paragraph{Finding 11: Random matching is liquidity-inefficient.}
Average liquidity of 52.2\% means that nearly half of all matched pairs fail to trade.  With random matching, many pairs are matched across non-overlapping value/cost ranges (buyer value $<$ seller cost), making trade impossible regardless of agent capability.

\paragraph{Single-seed caveat.}
These results are from a single seed (42).  The specific agent pairings, ZOPA distributions, and tick-level outcomes may vary under different seeds.  The findings are suggestive but should not be treated as confirmatory.

% ───────────────────────────────────────────────────────────────────
\subsection{Experiment E: Shock Response}
\label{sec:results-shock}

Table~\ref{tab:shock} compares market outcomes under the no-shock baseline and the shock condition.

\begin{table}[ht]
\centering
\caption{Shock response outcomes: \texttt{no\_shock} vs.\ \texttt{with\_shock} (30 ticks each, single seed).}
\label{tab:shock}
\begin{tabular}{@{}lccc@{}}
\toprule
Metric & \texttt{no\_shock} & \texttt{with\_shock} & $\Delta$ \\
\midrule
Mean price              & \$90.92 & \$88.77 & $-$\$2.15 \\
Price min               & \$78.07 & \$69.95 & $-$\$8.12 \\
Price max               & \$104.76 & \$100.82 & $-$\$3.94 \\
Within-tick price std   & \$20.36 & \$19.55 & $-$\$0.81 \\
Mean liquidity          & 54.2\% & 51.1\% & $-$3.1\,pp \\
Liquidity min           & 26.7\% & 13.3\% & $-$13.4\,pp \\
Liquidity max           & 86.7\% & 86.7\% & 0 \\
\bottomrule
\end{tabular}
\end{table}

% TODO: Generate Figure — overlaid price time series (shock vs no-shock).
% \begin{figure}[ht]
% \centering
% \includegraphics[width=0.8\textwidth]{figures/shock_price_comparison.pdf}
% \caption{Mean transaction price per tick under no-shock and shock conditions (30 ticks each, single seed).}
% \label{fig:shock-price}
% \end{figure}

\paragraph{Finding 12: Shocks reduce prices and liquidity.}
The shock condition produces a lower mean price (\$88.77 vs.\ \$90.92) and lower average liquidity (51.1\% vs.\ 54.2\%).  Random demand and supply multipliers in $[0.7, 1.3]$ with 30\% probability compress the effective ZOPA for affected pairs, reducing both the number of successful trades and the price level at which they occur.  H-E2 is confirmed.

\paragraph{Finding 13: Shocks produce asymmetric liquidity tail risk.}
The worst-case tick liquidity drops sharply with shocks (13.3\% vs.\ 26.7\%), while the best-case tick remains unchanged (86.7\% in both conditions).  Negative shocks that simultaneously suppress buyer values and inflate seller costs can eliminate nearly all viable ZOPA pairs in a single tick.  Positive shocks widen the ZOPA but cannot push liquidity above the random-matching ceiling.  H-E3 is confirmed.

\paragraph{Finding 14: Price dispersion is unchanged by shocks.}
Within-tick price standard deviation is nearly identical across conditions (\$20.36 vs.\ \$19.55).  The bilateral bargaining protocol---in which each pair negotiates independently---is the primary driver of price dispersion, not the width of the ZOPA.  The shock mechanism shifts the ZOPA but does not change the negotiation process.  H-E1 is partially rejected: while shocks affect price levels and liquidity, they do not increase within-tick price dispersion.

\paragraph{Single-seed caveat.}
As with Experiment~D, these results are from a single seed.  The specific shock realisations (which ticks are shocked, by how much) are fixed; different seeds would produce different shock sequences and potentially different aggregate effects.

% ───────────────────────────────────────────────────────────────────
\subsection{Experiment H: Institutional Mechanism Comparison}
\label{sec:results-mechanism}

Table~\ref{tab:mechanism} compares market outcomes under two matching mechanisms: random pairing and surplus-maximising (greedy maximum-ZOPA) pairing.

\begin{table}[ht]
\centering
\caption{Mechanism comparison outcomes (10 ticks $\times$ 10 pairs $\times$ 3 seeds per condition).}
\label{tab:mechanism}
\begin{tabular}{@{}lcc@{}}
\toprule
Metric & \texttt{random} & \texttt{surplus\_max} \\
\midrule
Total sessions          & 300   & 220 \\
Deals made              & 167   & 158 \\
Deal rate               & 55.7\% & 71.8\% \\
Mean price              & \$91.47 & \$87.24 \\
Price std               & \$18.60 & \$16.45 \\
Allocative efficiency   & 83.9\% & 89.7\% \\
Surplus Gini            & 0.226 & 0.216 \\
Buyer surplus (mean)    & \$30.54 & \$42.25 \\
Seller surplus (mean)   & \$17.39 & \$16.17 \\
\bottomrule
\end{tabular}
\end{table}

The surplus-maximising matcher produces fewer total sessions (220 vs.\ 300) because it only pairs agents with positive ZOPA overlap, leaving unprofitable pairs unmatched.  Despite fewer sessions, it achieves nearly as many deals (158 vs.\ 167) and a substantially higher deal rate (71.8\% vs.\ 55.7\%).

\paragraph{Finding 15: Surplus-maximising matching improves deal rates by 16 percentage points.}
The \texttt{surplus\_max} condition's deal rate of 71.8\% exceeds the \texttt{random} condition's 55.7\%.  By filtering out pairs with non-overlapping or narrow ZOPAs before negotiation begins, the mechanism ensures that matched pairs have sufficient surplus for the LLM agents to reach agreement.  H-H1 is confirmed.

\paragraph{Finding 16: Mechanism design improves allocative efficiency.}
Allocative efficiency rises from 83.9\% under random matching to 89.7\% under surplus-maximising matching, a 5.8 percentage point improvement.  The mechanism captures more of the theoretically available welfare by directing trades toward high-surplus pairs.  H-H2 is confirmed.

\paragraph{Finding 17: Buyer surplus increases disproportionately under surplus-maximising matching.}
Mean buyer surplus rises from \$30.54 to \$42.25 (+38\%), while seller surplus is essentially unchanged (\$17.39 vs.\ \$16.17).  Under surplus-maximising pairing, the additional surplus available from wider ZOPAs accrues predominantly to buyers, consistent with the first-mover advantage observed in Experiment~D.

\paragraph{Finding 18: Surplus-maximising matching reduces price dispersion.}
Price standard deviation falls from \$18.60 to \$16.45.  By excluding low-overlap pairs that produce extreme prices, the mechanism narrows the price distribution, though within-session dispersion remains substantial.

% ───────────────────────────────────────────────────────────────────
\subsection{Cross-Experiment Analysis}
\label{sec:results-cross}

Tables~\ref{tab:positive-results} and~\ref{tab:departures} synthesise the empirical patterns observed across all six experiments.

\begin{table}[ht]
\centering
\caption{Where LLM agents align with economic theory.}
\label{tab:positive-results}
\small
\begin{tabular}{@{}llc@{}}
\toprule
Economic principle & Evidence & Exp \\
\midrule
ZOPA determines deal success & 2.7\% vs.\ 100\% deal rate as ZOPA narrows & B \\
Market price convergence & +\$0.057/tick slope; stable \$89 equilibrium & D \\
Bilateral bargaining dispersion & Within-tick $\sigma = \$20$; Law of One Price fails & D \\
First-mover buyer advantage & Buyer surplus \$32.95 $>$ seller surplus \$16.40 & D \\
Shock $\to$ price/liquidity reduction & $-$\$2.15 price, $-$3.1\,pp liquidity & E \\
Asymmetric shock tail risk & Worst-case liquidity drops 13.4\,pp; best-case unchanged & E \\
Mechanism $\to$ efficiency gain & +5.8\,pp allocative efficiency under surplus-max matching & H \\
Mechanism $\to$ deal rate gain & +16.1\,pp deal rate under surplus-max matching & H \\
\bottomrule
\end{tabular}
\end{table}

\begin{table}[ht]
\centering
\caption{Where LLM agents depart from classical predictions.}
\label{tab:departures}
\small
\begin{tabular}{@{}lllc@{}}
\toprule
Predicted behaviour & Observed & Interpretation & Exp \\
\midrule
Anchoring bias & $r = -0.148$ & Prompt provides explicit surplus info & B \\
Deadline concession & Closes at round 3 always & Gate + prefer-counter steering & C \\
Shocks increase dispersion & $\sigma$ unchanged & Protocol drives dispersion, not ZOPA & E \\
\bottomrule
\end{tabular}
\end{table}

\paragraph{Feasibility gate as cross-cutting confound.}
The gate drives approximately 40\% of all acceptances.  In fixed-scenario experiments (A--C), it fires at round~2 in every session, limiting concession data to two observations and rendering settlement timing invariant to deadline manipulation.  In market experiments (D--E), gate involvement varies by session depending on the drawn ZOPA.  The gate's contribution to deal rates and settlement timing cannot be separated from agent capability without ablation studies that disable the gate entirely.

\paragraph{First-mover advantage as protocol effect.}
Buyer surplus exceeds seller surplus in both the fixed-scenario concession experiment (A) and the market experiment (D).  This pattern is consistent with the buyer-first turn order in the alternating-offer protocol: the buyer's opening offer anchors the negotiation toward the lower end of the ZOPA.  This is a protocol effect, not an agent-type effect, and would persist regardless of the LLM used.

\paragraph{Prompt steering as uncontrolled variable.}
The prefer-counter instruction and the accept-condition block are active in all LLM conditions across all experiments.  Their combined effect is to reduce rejection, accelerate settlement, and possibly suppress anchoring.  Because neither has been ablated, they function as uncontrolled confounds rather than experimental variables.  Separating their effects from the model's intrinsic bargaining behaviour is an important direction for future work.

% TODO: Consider adding a paragraph on chain-of-thought effect (deliberative vs reactive) as a cross-cutting finding if more data becomes available from gate-disabled runs.
